\begin{abstract}
\noindent 

The design of social security systems has implications %not only for households' saving and labor supply choices, but also
for the political support of intergenerational transfers. 
We examine the effects of contributive vs.\ universal pensions, as well as the link between benefits and earnings, on the social security tax rate, labor supply, and capital accumulation. Employing an overlapping generations model featuring heterogeneous households, our analysis aligns with existing results on the effects of retirees' relative political power, aging, and income inequality. %The analysis employs an overlapping generations model with households that differ in their labor productivity. Comparative statics are consistent with the existing literature concerning the effects of relative political power of retirees, ageing or income inequality. 
Additionally we find taxes to be higher in more Beveridgean systems and in more universal ones, but only when the share of households benefiting from universal pensions is large. %, with endogenous system characteristics that respond to demographic shifts, income inequality, and the share of displaced workers. %We show that the economic equilibrium features both lower formal labor supply and private savings, when pension benefits are universal. 
%We develop a recursive numerical approach that solves for the politico-economic equilibrium policy not only in steady state, but also as a function of time varying demographics.
%Due to a universal system's better insurance properties, society always prefers it over a contributive alternative. Nevertheless, higher tax rates might lead to lower steady state welfare under a universal system. %We show that an endogenous-gridpoint method is by far the computationally superior approach. 
%We conduct two numerical evaluations of the model's performance. First, 
%We calibrate the model to Argentina, which between 2005 and 2010 introduced a number of reforms to make coverage universal. W
When calibrated, the model explains the observed increase in pension spending of about 1.9\% of GDP following Argentina's 2005-2010 reforms to make coverage more universal. The model predicts that in rich countries the main determinants of pension spending are population ageing and the link between benefits and earnings. In contrast, inequality in income, or voting patterns, and leisure preferences have a minor effect.%if the shift in leisure preferences following the COVID-19 pandemic in the U.S. were to persist, this would result in a 1.8 p.p.\ increase in the social security tax rate.


\vspace{.3cm}

\noindent 
JEL Classification: D72, H55, J22, O17

\vspace{.3cm}
\noindent Keywords: Social Security, Labor supply, Politico-economic equilibrium, Heterogeneous agents.
\end{abstract}